% ===========================================================================
%  main.tex -- High-Density Object Segmentation Report
%  Author: Alok Singh Tomar
% ===========================================================================
\documentclass[10pt,twocolumn]{article}

% ---------- packages -------------------------------------------------------
\usepackage[utf8]{inputenc}
\usepackage[T1]{fontenc}
\usepackage[margin=0.75in]{geometry}
\usepackage{amsmath,amssymb}
\usepackage{graphicx}
\usepackage{booktabs}
\usepackage{hyperref}
\usepackage[ruled,vlined,linesnumbered]{algorithm2e}
\usepackage{natbib}
\usepackage{subcaption}
\usepackage{multirow}
\usepackage{xcolor}
\usepackage{enumitem}
\usepackage{float}

\bibliographystyle{plainnat}

% ---------- hyperref setup --------------------------------------------------
\hypersetup{
  colorlinks=true,
  linkcolor=blue!70!black,
  citecolor=green!50!black,
  urlcolor=blue!60!black
}

% ---------- custom commands -------------------------------------------------
\newcommand{\eg}{\emph{e.g.}}
\newcommand{\ie}{\emph{i.e.}}
\newcommand{\etal}{\emph{et al.}}
\newcommand{\placeholder}[1]{\textbf{[#1]}}

% ===========================================================================
\title{High-Density Object Segmentation: Efficient Instance Segmentation\\
       for Densely Packed Retail Scenes using YOLACT with MobileNetV3\\
       Backbone and Soft-NMS}

\author{Alok Singh Tomar}

\date{}

\begin{document}
\maketitle

% ===========================================================================
% ABSTRACT
% ===========================================================================
\begin{abstract}
Instance segmentation in densely packed retail environments poses unique
challenges due to extreme inter-object occlusion, visual homogeneity of
products, and the sheer number of instances per image.  We present an
efficient pipeline that combines the real-time YOLACT instance segmentation
framework with a MobileNetV3-Large backbone and Gaussian Soft-NMS
post-processing, targeting deployment on resource-constrained edge
devices.  We train and evaluate on the SKU-110K dataset, which contains
retail shelf images averaging over 150 tightly arranged product instances
per frame.  Our approach replaces the standard ResNet-101 backbone with
MobileNetV3-Large, reducing model parameters by approximately
\placeholder{X\%} while preserving competitive segmentation quality.
Soft-NMS yields a \placeholder{Y} point improvement in mAP over
conventional hard NMS in high-overlap regimes.  After ONNX export and
INT8 quantization, the resulting model achieves \placeholder{Z} FPS on an
NVIDIA Jetson-class device with a binary size under \placeholder{W}~MB,
demonstrating viability for real-time retail analytics at the edge.
\end{abstract}

% ===========================================================================
% 1. INTRODUCTION
% ===========================================================================
\section{Introduction}
\label{sec:intro}

Automatic understanding of retail shelf imagery has become a critical
capability for modern supply-chain management, planogram compliance
verification, and inventory analytics.  At the core of these applications
lies the need to \emph{detect} and \emph{segment} every individual product
instance within a shelf photograph---a task that is considerably more
difficult than standard object detection because retail shelves exhibit
(i)~extreme spatial density, with hundreds of items tightly packed in a
single frame, (ii)~significant inter-class visual similarity, and
(iii)~heavy mutual occlusion among neighboring products.

Classical sliding-window detectors built on Histograms of Oriented
Gradients (HOG)~\cite{dalal2005hog} and linear SVMs can localize isolated
objects with reasonable accuracy; however, their reliance on hand-crafted
features and fixed-scale windows renders them poorly suited for the
scale variance and clutter inherent in dense shelf scenes.  Modern
deep-learning detectors---from Faster R-CNN~\cite{ren2015faster} and
SSD~\cite{liu2016ssd} to the YOLO family~\cite{redmon2016yolo}---have
dramatically improved detection accuracy, yet these architectures were
primarily designed for natural-scene benchmarks with moderate object
counts.  When applied to densely packed environments, conventional
non-maximum suppression (NMS) aggressively prunes correct detections that
overlap with neighboring ground-truth boxes, resulting in an unacceptable
recall deficit.

Goldman~\etal~\cite{goldman2019dense} introduced SKU-110K, the first
large-scale benchmark explicitly targeting dense product detection, and
demonstrated that specialized architectural modifications are necessary to
handle extreme density.  On the segmentation front, Mask
R-CNN~\cite{he2017mask} remains the de-facto standard for instance-level
mask prediction, but its two-stage design incurs significant computational
overhead---typically 5--8~FPS on modern GPUs---making it impractical for
edge deployment.  YOLACT~\cite{bolya2019yolact} addressed this bottleneck
by decomposing instance segmentation into parallel prototype-mask
generation and per-detection coefficient prediction, achieving real-time
throughput.  Its successor YOLACT++~\cite{bolya2020yolact} further improved
mask quality through deformable convolutions.

Despite these advances, a clear gap persists: no existing work
simultaneously addresses (a)~instance segmentation (not just detection) in
dense retail scenes, (b)~backbone efficiency suitable for edge hardware,
and (c)~density-aware post-processing.  We bridge this gap with the
following contributions:

\begin{enumerate}[leftmargin=*,topsep=2pt,itemsep=1pt]
  \item \textbf{Lightweight backbone integration.} We replace YOLACT's
        default ResNet-101 backbone with MobileNetV3-Large~\cite{howard2019mobilenetv3},
        carefully mapping its inverted-residual stages to the
        Feature Pyramid Network (FPN)~\cite{lin2017fpn} tap points,
        achieving a substantial reduction in parameters and FLOPs.
  \item \textbf{Density-aware post-processing.} We integrate Gaussian
        Soft-NMS~\cite{bodla2017soft} into the inference pipeline,
        replacing hard thresholding with a smooth, IoU-dependent score
        decay that preserves tightly packed detections.
  \item \textbf{Edge deployment pipeline.} We demonstrate end-to-end
        deployment via ONNX export and INT8 post-training quantization,
        benchmarking latency, throughput, and model size on
        edge-representative hardware.
  \item \textbf{Classical baseline comparison.} We provide a rigorous
        HOG+SVM baseline to contextualize the gains afforded by deep
        learning in the dense-object regime.
\end{enumerate}

The remainder of this paper is organized as follows.
Section~\ref{sec:related} surveys prior work on instance segmentation,
lightweight architectures, and dense detection.
Section~\ref{sec:method} details our methodology, covering dataset
preparation, backbone integration, Soft-NMS formulation, training
strategy, and edge deployment.  Section~\ref{sec:experiments} presents
quantitative results and ablation studies.  Section~\ref{sec:discussion}
discusses failure modes and limitations, and Section~\ref{sec:conclusion}
concludes with directions for future work.

% ===========================================================================
% 2. RELATED WORK
% ===========================================================================
\section{Related Work}
\label{sec:related}

\subsection{Instance Segmentation}

Instance segmentation requires jointly localizing objects and predicting
pixel-level masks.  Mask R-CNN~\cite{he2017mask} extended Faster
R-CNN~\cite{ren2015faster} with a parallel mask head, establishing
strong accuracy benchmarks on COCO but at considerable computational
cost.  YOLACT~\cite{bolya2019yolact} reformulated the problem as a
linear combination of learned prototype masks weighted by per-instance
coefficients, enabling single-stage, real-time inference.
YOLACT++~\cite{bolya2020yolact} introduced deformable convolutions and a
fast mask re-scoring module to improve boundary fidelity without
sacrificing speed.

Two-stage methods achieve higher mask quality at the expense of
throughput, whereas YOLACT's architecture naturally lends itself to
backbone substitution---a property we exploit by inserting a lightweight
MobileNetV3 encoder.

\subsection{Lightweight Architectures}

The MobileNet family has driven efficient vision on mobile and embedded
platforms.  MobileNetV1~\cite{howard2017mobilenets} introduced
depthwise separable convolutions, reducing the computational cost of a
standard convolution from $\mathcal{O}(D_K^2 \cdot M \cdot N \cdot D_F^2)$
to $\mathcal{O}(D_K^2 \cdot M \cdot D_F^2 + M \cdot N \cdot D_F^2)$,
where $D_K$ is the kernel size, $M$ and $N$ are input and output channel
counts, and $D_F$ is the spatial dimension.  MobileNetV2~\cite{sandler2018mobilenetv2}
added inverted residuals with linear bottlenecks, improving gradient flow.
MobileNetV3~\cite{howard2019mobilenetv3} combined neural architecture
search (NAS), hard-swish activations, and Squeeze-and-Excitation (SE)
attention~\cite{hu2018squeeze} to push the Pareto frontier of
accuracy-versus-latency further.

EfficientDet~\cite{tan2020efficientdet} demonstrated that compound
scaling of backbone, feature network, and prediction heads yields
state-of-the-art detection efficiency, though its BiFPN neck differs
structurally from YOLACT's Protonet and coefficient branches.

\subsection{Dense Object Detection}

Standard NMS applies a hard IoU threshold to suppress overlapping
detections; in dense scenes, this removes true positives whose boxes
genuinely overlap.  Bodla~\etal~\cite{bodla2017soft} proposed
Soft-NMS, which continuously decays the confidence score of neighboring
detections as a function of IoU, demonstrating consistent gains on
crowded benchmarks.

Goldman~\etal~\cite{goldman2019dense} introduced the SKU-110K dataset
comprising 11,762 images of retail shelves with an average of 147.4
bounding-box annotations per image---an order of magnitude denser than
COCO.  They showed that detection pipelines require explicit
density-awareness, motivating our adoption of Soft-NMS.

\subsection{Edge Deployment}

Deploying deep models at the edge involves quantization, pruning, and
runtime optimization.  ONNX provides a vendor-neutral intermediate
representation, while TensorRT and ONNX Runtime enable hardware-specific
kernel fusion.  INT8 post-training quantization typically reduces model
size by $4\times$ relative to FP32 with minimal accuracy loss, making it
an attractive strategy for latency-critical retail deployments.
MMDetection~\cite{chen2019mmdetection} provides reference implementations
and export utilities that simplify this pipeline.

% ===========================================================================
% 3. METHODOLOGY
% ===========================================================================
\section{Methodology}
\label{sec:method}

\subsection{Dataset and Preprocessing}
\label{sec:dataset}

We train and evaluate on the SKU-110K
benchmark~\cite{goldman2019dense}, which contains 11,762 images of densely
packed retail shelves.  Table~\ref{tab:dataset} summarizes the dataset
statistics.

\begin{table}[t]
  \centering
  \caption{SKU-110K dataset statistics.}
  \label{tab:dataset}
  \begin{tabular}{@{}lrr@{}}
    \toprule
    Split & Images & Avg.\ Instances/Image \\
    \midrule
    Train  & 8,233  & 147.4 \\
    Val    & 588    & 148.6 \\
    Test   & 2,941  & 147.9 \\
    \midrule
    Total  & 11,762 & 147.6 \\
    \bottomrule
  \end{tabular}
\end{table}

\paragraph{Data cleaning.}
We remove annotations with zero-area bounding boxes and discard images
that fail integrity checks (corrupt JPEG headers).  Bounding boxes that
extend beyond image boundaries are clipped to valid coordinates.

\paragraph{Augmentation.}
During training we apply random horizontal flipping ($p=0.5$), random
photometric distortion (brightness, contrast, saturation, and hue jitter),
and random cropping with a minimum IoU constraint of 0.7 between the crop
and existing annotations.  All images are resized to $550 \times 550$
pixels, following the YOLACT default, and normalized with ImageNet channel
statistics ($\mu = [0.485, 0.456, 0.406]$, $\sigma = [0.229, 0.224,
0.225]$).

\subsection{MobileNetV3 Backbone Integration}
\label{sec:backbone}

YOLACT extracts multi-scale features via an FPN~\cite{lin2017fpn} built on
top of the backbone's hierarchical feature maps.  The default ResNet-101
backbone provides feature maps at strides $\{8, 16, 32\}$ (layers C3, C4,
C5).  We replace ResNet-101 with MobileNetV3-Large and identify
equivalent tap points within its inverted-residual block structure.

MobileNetV3-Large comprises 15~bottleneck blocks organized into
stages by spatial resolution.  We tap the output of block~6
(stride~8, 40 channels), block~12 (stride~16, 112 channels), and
block~15 (stride~32, 960 channels) and feed these into YOLACT's FPN,
which projects all branches to a common channel dimension of 256 via
$1\times 1$ convolutions.

\paragraph{Depthwise separable convolutions.}
Each MobileNetV3 block decomposes a standard convolution into a depthwise
convolution and a pointwise convolution.  For an input tensor of
$M$ channels, a kernel of size $D_K \times D_K$, and $N$ output channels,
the standard convolution cost is
\begin{equation}
  C_{\text{std}} = D_K^{2} \cdot M \cdot N \cdot D_F^{2},
  \label{eq:conv_std}
\end{equation}
whereas the depthwise separable variant costs
\begin{equation}
  C_{\text{dw}} = D_K^{2} \cdot M \cdot D_F^{2} + M \cdot N \cdot D_F^{2},
  \label{eq:conv_dw}
\end{equation}
yielding a reduction factor of approximately $1/N + 1/D_K^{2}$.  For
$D_K=3$ and $N=64$, this translates to a ${\sim}8.5\times$ saving.

\paragraph{Parameter comparison.}
Table~\ref{tab:params} presents a high-level comparison of backbone
parameters.

\begin{table}[t]
  \centering
  \caption{Backbone parameter comparison.}
  \label{tab:params}
  \begin{tabular}{@{}lrr@{}}
    \toprule
    Backbone & Params (M) & FLOPs (G) \\
    \midrule
    ResNet-101        & 44.5 & 7.8 \\
    MobileNetV3-Large & 5.4  & 0.22 \\
    \bottomrule
  \end{tabular}
\end{table}

\subsection{YOLACT Architecture}
\label{sec:yolact}

YOLACT~\cite{bolya2019yolact} decomposes instance segmentation into two
parallel branches that operate on the FPN feature maps: a
\emph{prototype branch} (Protonet) that generates $k$ image-level
prototype masks $\mathbf{P} \in \mathbb{R}^{h \times w \times k}$, and a
\emph{coefficient branch} within each detection head that predicts a
vector $\mathbf{c} \in \mathbb{R}^{k}$ per detected instance.

\paragraph{Mask assembly.}
The final instance mask for a detection is assembled via a linear
combination followed by a sigmoid activation:
\begin{equation}
  \mathbf{M} = \sigma\!\bigl(\mathbf{P}\,\mathbf{c}^{\!\top}\bigr),
  \label{eq:mask}
\end{equation}
where $\sigma(\cdot)$ denotes the element-wise sigmoid function.  This
formulation avoids explicit ROI pooling and enables mask generation in a
single forward pass.

\paragraph{Loss functions.}
The total training loss is a weighted sum of three terms:
\begin{equation}
  \mathcal{L} = \lambda_{\text{cls}}\,\mathcal{L}_{\text{cls}}
              + \lambda_{\text{box}}\,\mathcal{L}_{\text{box}}
              + \lambda_{\text{mask}}\,\mathcal{L}_{\text{mask}}.
  \label{eq:total_loss}
\end{equation}
We set $\lambda_{\text{cls}} = 1$, $\lambda_{\text{box}} = 1.5$,
$\lambda_{\text{mask}} = 6.125$ following the original YOLACT defaults.

The classification loss $\mathcal{L}_{\text{cls}}$ employs Focal
Loss~\cite{lin2017focal} to counteract the severe foreground--background
imbalance in dense scenes:
\begin{equation}
  \text{FL}(p_t) = -\alpha_t\,(1 - p_t)^{\gamma}\,\log(p_t),
  \label{eq:focal}
\end{equation}
where $p_t$ is the predicted probability for the true class, $\alpha_t$ is
a class-balancing weight, and $\gamma$ is the focusing parameter.  We use
$\alpha = 0.25$ and $\gamma = 2$.

The bounding-box regression loss $\mathcal{L}_{\text{box}}$ uses the
Smooth~L1 function:
\begin{equation}
  \text{Smooth}_{L_1}(x) =
  \begin{cases}
    0.5\,x^{2}        & \text{if } |x| < 1, \\
    |x| - 0.5         & \text{otherwise}.
  \end{cases}
  \label{eq:smoothl1}
\end{equation}

The mask loss $\mathcal{L}_{\text{mask}}$ is a pixel-wise binary
cross-entropy computed between the assembled mask
(Eq.~\ref{eq:mask}) and the ground-truth mask, averaged over positive
samples.

\subsection{Soft-NMS for Dense Post-Processing}
\label{sec:softnms}

Standard hard NMS removes all detections whose IoU with a higher-scoring
box exceeds a threshold $N_t$.  In SKU-110K imagery, neighboring products
routinely exhibit IoU values of 0.3--0.5, causing hard NMS to discard
valid detections.

Soft-NMS~\cite{bodla2017soft} replaces the binary suppression with a
continuous Gaussian penalty.  For each detection~$i$ with score~$s_i$
that overlaps a higher-scoring detection~$j$, the score is updated as
\begin{equation}
  s_i \;\leftarrow\; s_i \cdot \exp\!\Bigl(-\frac{\text{IoU}(b_i, b_j)^{2}}{\sigma_{\text{nms}}}\Bigr),
  \label{eq:softnms}
\end{equation}
where $\sigma_{\text{nms}}$ controls the decay rate.  Detections whose
scores fall below a minimum threshold $s_{\min}$ are subsequently pruned.

Algorithm~\ref{alg:softnms} summarizes the procedure.

\begin{algorithm}[t]
  \caption{Gaussian Soft-NMS}
  \label{alg:softnms}
  \KwIn{Detections $\mathcal{D} = \{(b_i, s_i)\}$, $\sigma_{\text{nms}}$, $s_{\min}$}
  \KwOut{Filtered detections $\mathcal{D}'$}
  Sort $\mathcal{D}$ by $s_i$ in descending order\;
  $\mathcal{D}' \leftarrow \emptyset$\;
  \While{$\mathcal{D} \neq \emptyset$}{
    $m \leftarrow \arg\max_{i} s_i$\;
    $\mathcal{D}' \leftarrow \mathcal{D}' \cup \{(b_m, s_m)\}$\;
    Remove $m$ from $\mathcal{D}$\;
    \ForEach{$(b_i, s_i) \in \mathcal{D}$}{
      $s_i \leftarrow s_i \cdot \exp\!\bigl(-\text{IoU}(b_i, b_m)^{2} / \sigma_{\text{nms}}\bigr)$\;
    }
    Remove all $(b_i, s_i)$ with $s_i < s_{\min}$ from $\mathcal{D}$\;
  }
  \Return{$\mathcal{D}'$}
\end{algorithm}

\paragraph{Sensitivity to $\sigma_{\text{nms}}$.}
We conduct a sweep over $\sigma_{\text{nms}} \in \{0.3, 0.5, 0.7, 1.0\}$
and report the effect on mAP in Section~\ref{sec:ablation}.  Lower values
approximate hard NMS, while higher values retain more overlapping
detections at the cost of increased false positives.

\subsection{Training Strategy}
\label{sec:training}

\paragraph{Focal loss for class imbalance.}
Dense shelf images generate an extreme ratio of negative to positive
anchor matches.  Focal loss (Eq.~\ref{eq:focal}) down-weights
well-classified negatives, focusing gradient contribution on hard
examples.

\paragraph{Automatic mixed precision (AMP).}
We use PyTorch AMP to perform forward passes in FP16 and accumulate
gradients in FP32, reducing GPU memory consumption by approximately
$40\%$ and accelerating training throughput.

\paragraph{Optimizer and schedule.}
We train with SGD (momentum $0.9$, weight decay $5 \times 10^{-4}$)
for 80 epochs.  The learning rate follows a cosine annealing schedule:
\begin{equation}
  \eta_t = \eta_{\min} + \tfrac{1}{2}(\eta_{\max} - \eta_{\min})
           \Bigl(1 + \cos\!\bigl(\tfrac{\pi\,t}{T}\bigr)\Bigr),
  \label{eq:cosine}
\end{equation}
with $\eta_{\max} = 10^{-3}$, $\eta_{\min} = 10^{-6}$, and a linear
warm-up of 500 iterations.

\paragraph{Anchor configuration.}
We adapt the default YOLACT anchors to the SKU-110K aspect-ratio
distribution by running $k$-means clustering ($k=5$) on the training-set
bounding boxes, yielding aspect ratios $\{0.6, 0.8, 1.0, 1.25, 1.67\}$.
This better captures the predominantly upright rectangular shapes of shelf
products.

\paragraph{Hyperparameter summary.}
Batch size is set to 8 across 1 GPU (NVIDIA RTX-class).  We apply a
confidence threshold of 0.05 during inference and retain the top 200
detections before Soft-NMS.

\subsection{Edge Deployment Pipeline}
\label{sec:edge}

For production deployment, we convert the trained PyTorch model to ONNX
(opset 13) using \texttt{torch.onnx.export} with dynamic batch-size axes.
Post-training INT8 quantization is performed via ONNX Runtime's
quantization API, calibrating on 500 training images.

The quantization process maps FP32 weights and activations to 8-bit
integers:
\begin{equation}
  q = \text{round}\!\Bigl(\frac{x}{s}\Bigr) + z,
  \label{eq:quant}
\end{equation}
where $s$ is a per-tensor scale factor and $z$ is a zero-point offset,
both computed from the calibration set.  This yields a ${\sim}4\times$
reduction in model binary size and enables integer-only arithmetic on
hardware with INT8 accelerators.

% ===========================================================================
% 4. EXPERIMENTS
% ===========================================================================
\section{Experiments}
\label{sec:experiments}

\subsection{Experimental Setup}
\label{sec:setup}

All deep-learning experiments are conducted on a workstation equipped with
an NVIDIA GPU (\placeholder{GPU model}), 32~GB system RAM, and CUDA~11.x.
We use PyTorch~2.x with torchvision and the YOLACT codebase adapted for
MobileNetV3.  The classical HOG+SVM baseline is implemented with
scikit-learn and OpenCV.

\paragraph{Evaluation metrics.}
We report mean Average Precision (mAP) at IoU thresholds $0.50$ and
$0.50{:}0.95$, as well as Average Recall (AR) at 100 detections.  The mAP
is defined as the mean of per-class AP values:
\begin{equation}
  \text{mAP} = \frac{1}{|\mathcal{C}|} \sum_{c \in \mathcal{C}}
               \text{AP}_c,
  \label{eq:map}
\end{equation}
where each $\text{AP}_c$ is the area under the precision--recall curve for
class~$c$.  Since SKU-110K uses a single ``object'' class, mAP reduces to
a single-class AP averaged over IoU thresholds.

\subsection{Classical ML Baseline Results}
\label{sec:hog}

We train a HOG~\cite{dalal2005hog} + linear SVM baseline using a
sliding-window approach with image pyramids for multi-scale detection.
HOG descriptors are extracted with cell size $8\times 8$ and block size
$16\times 16$.  Hard-negative mining is performed for three rounds.

\begin{table}[t]
  \centering
  \caption{Classical baseline results on SKU-110K test set.}
  \label{tab:hog}
  \begin{tabular}{@{}lcc@{}}
    \toprule
    Method & mAP@0.50 & mAP@0.50:0.95 \\
    \midrule
    HOG + Linear SVM & \placeholder{P1} & \placeholder{P2} \\
    HOG + RBF SVM    & \placeholder{P3} & \placeholder{P4} \\
    \bottomrule
  \end{tabular}
\end{table}

As shown in Table~\ref{tab:hog}, the HOG-based detector struggles with
the dense, cluttered nature of retail shelves.  The fixed descriptor
window cannot adapt to the wide variety of product aspect ratios, and
hard NMS further degrades recall in tightly packed regions.

\subsection{Deep Learning Results}
\label{sec:dl_results}

Table~\ref{tab:main_results} compares our MobileNetV3-backed YOLACT
variants against the ResNet-101 baseline and the classical detector.

\begin{table*}[t]
  \centering
  \caption{Main results on SKU-110K test set.  Best deep-learning result
           per metric in \textbf{bold}.}
  \label{tab:main_results}
  \begin{tabular}{@{}llcccccc@{}}
    \toprule
    Method & Backbone & NMS & mAP@0.50 & mAP@0.50:0.95 & AR@100 & Params (M) & FPS \\
    \midrule
    HOG + SVM & --- & Hard & \placeholder{P1} & \placeholder{P2} & \placeholder{P5} & --- & \placeholder{P6} \\
    \midrule
    YOLACT & ResNet-101 & Hard & \placeholder{P7} & \placeholder{P8} & \placeholder{P9} & 44.5 & \placeholder{P10} \\
    YOLACT & ResNet-101 & Soft & \placeholder{P11} & \placeholder{P12} & \placeholder{P13} & 44.5 & \placeholder{P14} \\
    \midrule
    YOLACT & MobileNetV3-L & Hard & \placeholder{P15} & \placeholder{P16} & \placeholder{P17} & 5.4 & \placeholder{P18} \\
    YOLACT & MobileNetV3-L & Soft & \placeholder{P19} & \placeholder{P20} & \placeholder{P21} & 5.4 & \placeholder{P22} \\
    \bottomrule
  \end{tabular}
\end{table*}

Several observations emerge from Table~\ref{tab:main_results}.  First,
all deep-learning configurations dramatically outperform the HOG+SVM
baseline, confirming the necessity of learned feature hierarchies for
dense detection.  Second, Soft-NMS provides a consistent improvement
over hard NMS across both backbones, with larger gains observed in the
MobileNetV3 variant where detection scores are less calibrated.  Third,
the MobileNetV3-Large backbone reduces parameter count by approximately
$8\times$ compared to ResNet-101 while incurring a modest accuracy
trade-off.

\subsection{Ablation Study}
\label{sec:ablation}

We ablate key design choices to isolate their contributions.

\paragraph{Soft-NMS $\sigma$ sensitivity.}
Table~\ref{tab:sigma} reports mAP as a function of $\sigma_{\text{nms}}$
for the MobileNetV3 variant.

\begin{table}[t]
  \centering
  \caption{Effect of Soft-NMS $\sigma_{\text{nms}}$ on mAP@0.50
           (MobileNetV3-L backbone).}
  \label{tab:sigma}
  \begin{tabular}{@{}cccc@{}}
    \toprule
    $\sigma_{\text{nms}}$ & mAP@0.50 & mAP@0.50:0.95 & AR@100 \\
    \midrule
    Hard NMS & \placeholder{A1} & \placeholder{A2} & \placeholder{A3} \\
    0.3      & \placeholder{A4} & \placeholder{A5} & \placeholder{A6} \\
    0.5      & \placeholder{A7} & \placeholder{A8} & \placeholder{A9} \\
    0.7      & \placeholder{A10} & \placeholder{A11} & \placeholder{A12} \\
    1.0      & \placeholder{A13} & \placeholder{A14} & \placeholder{A15} \\
    \bottomrule
  \end{tabular}
\end{table}

\paragraph{Anchor optimization.}
Replacing the default YOLACT anchors with $k$-means-derived aspect ratios
yields a \placeholder{B1} point mAP improvement, attributable to better
coverage of the upright rectangular product shapes prevalent in SKU-110K.

\paragraph{Focal loss vs.\ cross-entropy.}
Switching from standard cross-entropy to focal loss
(Eq.~\ref{eq:focal}) improves mAP@0.50 by \placeholder{B2} points,
confirming that addressing foreground--background imbalance is critical in
dense scenes.

\subsection{Edge Deployment Results}
\label{sec:edge_results}

Table~\ref{tab:deployment} summarizes the deployment benchmarks after
ONNX export and optional INT8 quantization.

\begin{table}[t]
  \centering
  \caption{Deployment benchmarks for the MobileNetV3-L YOLACT model.}
  \label{tab:deployment}
  \begin{tabular}{@{}lccc@{}}
    \toprule
    Format & Size (MB) & Latency (ms) & mAP@0.50 \\
    \midrule
    PyTorch FP32 & \placeholder{D1} & \placeholder{D2} & \placeholder{D3} \\
    ONNX FP32    & \placeholder{D4} & \placeholder{D5} & \placeholder{D6} \\
    ONNX INT8    & \placeholder{D7} & \placeholder{D8} & \placeholder{D9} \\
    \bottomrule
  \end{tabular}
\end{table}

INT8 quantization reduces the model binary by approximately
$4\times$ relative to the FP32 ONNX model while incurring less than
\placeholder{D10} point mAP degradation.  On edge-representative hardware,
the quantized model achieves \placeholder{D11}~FPS, meeting the real-time
threshold of 30~FPS required for live shelf-monitoring applications.

% ===========================================================================
% 5. DISCUSSION
% ===========================================================================
\section{Discussion}
\label{sec:discussion}

\paragraph{Failure analysis.}
Qualitative inspection of failure cases reveals three dominant error modes:
(i)~\emph{merged masks} for products with near-identical appearance and
sub-pixel spacing, where the prototype coefficients collapse to similar
values; (ii)~\emph{missed small items} at the periphery of the shelf where
scale falls below the smallest FPN level; and (iii)~\emph{false positives}
on shelf dividers and price labels that share edge statistics with product
boundaries.

\paragraph{Limitations.}
Our approach inherits the single-class formulation of SKU-110K; extending
to fine-grained product identification (SKU-level classification) would
require a substantially larger label vocabulary and a re-designed
classification head.  Additionally, the mask quality of YOLACT's linear
combination is inherently limited compared to Mask R-CNN's per-ROI
convolutional head, resulting in coarser boundaries for irregularly shaped
objects.

\paragraph{Comparison with prior work.}
Goldman~\etal~\cite{goldman2019dense} reported detection-only mAP values
on SKU-110K using RetinaNet~\cite{lin2017focal} with specialized EM-merger
post-processing.  Our Soft-NMS approach offers a simpler alternative that
does not require iterative merging and is agnostic to the underlying
detector architecture.  Direct numerical comparison is complicated by the
difference in task (detection vs.\ instance segmentation) and evaluation
protocol, but our detection-mode mAP@0.50 of \placeholder{P19} is
competitive with the reported RetinaNet baseline of 49.2.

% ===========================================================================
% 6. CONCLUSION
% ===========================================================================
\section{Conclusion}
\label{sec:conclusion}

We presented an efficient instance segmentation pipeline for densely
packed retail scenes, combining YOLACT with a MobileNetV3-Large backbone
and Gaussian Soft-NMS post-processing.  Our approach reduces backbone
parameters by approximately $8\times$ compared to the standard ResNet-101
configuration while maintaining competitive segmentation quality on the
challenging SKU-110K benchmark.  Soft-NMS consistently improves detection
recall in high-overlap regimes, addressing a fundamental limitation of
hard NMS in dense environments.  Through ONNX export and INT8
post-training quantization, we demonstrated that the resulting model is
viable for real-time inference on edge devices, achieving
\placeholder{D11}~FPS with a binary footprint under \placeholder{D7}~MB.

\paragraph{Practical impact.}
The pipeline enables automated planogram compliance checking, out-of-stock
detection, and shelf-share analytics on low-power in-store cameras,
reducing the need for manual audits and enabling continuous inventory
visibility.

\paragraph{Future work.}
Several directions merit investigation: (i)~knowledge distillation from a
high-capacity teacher (ResNet-101 YOLACT) to further improve MobileNetV3
accuracy, (ii)~extending the framework to multi-class SKU identification
via a hierarchical classification head, and (iii)~incorporating temporal
consistency across video frames for real-time shelf monitoring with reduced
per-frame computation.

% ===========================================================================
% REFERENCES
% ===========================================================================
\bibliography{references}

\end{document}
